\ifdefined\BookMaster
\def\BookEnd{}
\else
\documentclass[a4paper,12pt,fleqn,openany,twoside,fontset=windows]{ctexbook}

\usepackage[fleqn]{amsmath}
\usepackage{amssymb,amsthm,mathrsfs}
\usepackage{geometry}
\usepackage[dvipsnames]{xcolor}
\usepackage{graphicx}
\usepackage{tikz}
\usetikzlibrary{calc}
\usepackage[most]{tcolorbox}
\usepackage{fancyhdr}
\usepackage{titlesec}
\usepackage{tocloft}
\usepackage{enumitem}
\usepackage{microtype}
\usepackage{pifont}
\usepackage{hyperref}
\usepackage{romannum}
\usepackage{mathtools}

\definecolor{bookblue}{HTML}{264653}
\definecolor{bookteal}{HTML}{4F7C82}
\definecolor{booklight}{HTML}{EDF4F3}
\definecolor{bookgray}{HTML}{5E686B}

\geometry{
    inner=2.7cm,
    outer=2.5cm,
    top=2.7cm,
    bottom=2.7cm,
    headheight=18pt,
    headsep=18pt,
    footskip=25pt
}
\setlength{\mathindent}{0pt}
\setlength{\parindent}{5pt}
\setlength{\parskip}{3pt plus 1pt minus 1pt}
\linespread{1.25}
\allowdisplaybreaks[3]
\AtBeginDocument{%
    \setlength{\parindent}{5pt}
    \setlength{\abovedisplayskip}{8pt plus 2pt minus 2pt}
    \setlength{\belowdisplayskip}{8pt plus 2pt minus 2pt}
    \setlength{\abovedisplayshortskip}{5pt plus 1pt minus 1pt}
    \setlength{\belowdisplayshortskip}{5pt plus 1pt minus 1pt}
}

\hypersetup{
    unicode=true,
    colorlinks=true,
    linkcolor=bookblue,
    citecolor=bookblue,
    urlcolor=bookteal,
    bookmarksopen=true,
    pdfauthor={Chen},
    pdftitle={数学分析习题整理}
}

\renewcommand{\chaptermark}[1]{%
    \edef\@tmp{第\thechapter 章\quad #1}%
    \expandafter\markboth\expandafter{\@tmp}{}%
}
\fancypagestyle{main}{%
    \fancyhf{}
    \fancyhead[LE]{\small\color{bookgray}\nouppercase{\leftmark}}
    \fancyhead[RO]{\small\color{bookgray}数学分析习题整理}
    \fancyfoot[C]{\color{bookblue}\thepage}
    \fancyfoot[R]{\small\bfseries Chen\;\; github.com/mathlover-1/math-analysis}
    \renewcommand{\headrulewidth}{0.4pt}
    \renewcommand{\footrulewidth}{0pt}
}
\fancypagestyle{plain}{%
    \fancyhf{}
    \fancyfoot[C]{\color{bookblue}\thepage}
    \fancyfoot[R]{\small\bfseries Chen\;\; github.com/mathlover-1/math-analysis}
    \renewcommand{\headrulewidth}{0pt}
    \renewcommand{\footrulewidth}{0pt}
}
\pagestyle{main}

\titleformat{\chapter}[display]
  {\normalfont\sffamily\bfseries\color{bookblue}}
  {\raggedleft\fontsize{46}{50}\selectfont\thechapter}
  {-18pt}
  {\titlerule[1.2pt]\vspace{10pt}\filright\Huge}
\titlespacing*{\chapter}{0pt}{-15pt}{36pt}

\renewcommand{\contentsname}{目\quad 录}
\renewcommand{\cfttoctitlefont}{\hfill\sffamily\bfseries\Huge\color{bookblue}}
\renewcommand{\cftaftertoctitle}{\hfill}
\setlength{\cftbeforetoctitleskip}{5pt}
\setlength{\cftaftertoctitleskip}{28pt}
\setlength{\cftbeforechapskip}{12pt}
\renewcommand{\cftchapfont}{\sffamily\bfseries\large\color{bookblue}}
\renewcommand{\cftchappagefont}{\sffamily\bfseries\large\color{bookblue}}
\renewcommand{\cftchapleader}{\color{bookteal}\cftdotfill{2.5}}

\newtcolorbox{ti}{
    enhanced,
    breakable,
    colback=booklight!55!white,
    colframe=bookteal!35!white,
    boxrule=0.45pt,
    boxsep=0pt,
    left=10pt,
    right=10pt,
    top=9pt,
    bottom=9pt,
    before skip=14pt,
    after skip=14pt,
    arc=2pt,
    outer arc=2pt,
    before upper={\parindent=0pt}
}

\newenvironment{booknotebody}{%
    \begin{center}
    \begin{minipage}{0.84\textwidth}
    \songti\zihao{-4}\linespread{1.45}\selectfont
    \setlength{\parindent}{2\ccwd}
    \setlength{\parskip}{0.75em}
}{%
    \end{minipage}
    \end{center}
}

\newenvironment{booknotelongbody}{%
    \begingroup
    \songti\zihao{-4}\linespread{1.45}\selectfont
    \setlength{\parindent}{2\ccwd}
    \setlength{\parskip}{0.75em}
    \begin{list}{}{%
        \setlength{\leftmargin}{0.08\textwidth}
        \setlength{\rightmargin}{0.08\textwidth}
        \setlength{\topsep}{0pt}
        \setlength{\partopsep}{0pt}
        \setlength{\parsep}{0pt}
        \setlength{\itemsep}{0pt}
    }
    \item\relax
}{%
    \end{list}
    \endgroup
}

\fancypagestyle{plainnowm}{%
    \fancyhf{}
    \fancyfoot[C]{\color{bookblue}\thepage}
    \renewcommand{\headrulewidth}{0pt}
    \renewcommand{\footrulewidth}{0pt}
}

\newcommand{\booknotetitle}[2][plain]{%
    \clearpage
    \phantomsection
    \addcontentsline{toc}{chapter}{#2}
    \markboth{#2}{}
    \thispagestyle{#1}
    \vspace*{1.2cm}
    \begin{center}
        {\sffamily\heiti\bfseries\fontsize{26}{34}\selectfont\color{bookblue}#2\par}
        \vspace{0.8em}
        {\color{bookteal}\rule{4.5cm}{0.8pt}\par}
    \end{center}
    \vspace{0.6cm}
}

\renewcommand{\proofname}{证明}
\renewcommand{\qedsymbol}{\color{bookteal}\ensuremath{\blacksquare}}

\title{数学分析习题整理}
\author{Chen}
\date{}

\makeatletter
% ==================== 旧版封面/封底设计备份 ====================
% 2026-08 重设计,旧代码用 \iffalse...\fi 封存,新版见下方。
\iffalse
\renewcommand{\maketitle}{%
    \begin{titlepage}
        \thispagestyle{empty}
        \setcounter{page}{0}
        \noindent\hspace*{-\parindent}\begin{tikzpicture}[remember picture,overlay]
            \fill[bookblue] (current page.north west) rectangle ([yshift=-4.2cm]current page.north east);
            \fill[booklight] ([yshift=-4.2cm]current page.north west) rectangle (current page.south east);
            \draw[bookteal!55,line width=0.7pt]
                ([xshift=2.4cm,yshift=4.3cm]current page.south west)
                .. controls +(2.5cm,1.8cm) and +(-2.8cm,-1.3cm) ..
                ([xshift=-2.4cm,yshift=7.2cm]current page.south east);
            \draw[bookteal!35,line width=0.7pt]
                ([xshift=2.5cm,yshift=6.2cm]current page.south west)
                .. controls +(3.0cm,-1.2cm) and +(-3.0cm,1.6cm) ..
                ([xshift=-2.3cm,yshift=3.8cm]current page.south east);
        \end{tikzpicture}
        \vspace*{5.6cm}
        \noindent\color{bookblue}\rule{\textwidth}{1.4pt}\par
        \vspace{8pt}
        {\centering\sffamily\bfseries\fontsize{32}{42}\selectfont\color{bookblue}\@title\par}
        \vspace{8pt}
        \noindent\color{bookteal}\rule{\textwidth}{0.6pt}\par
        \vspace{2.6cm}
        {\centering\Large\color{bookgray}\@author\par}
        \vfill
        {\centering\color{bookteal}\large
            $\displaystyle \lim_{n\to\infty} a_n$\qquad
            $\displaystyle \int_a^b f(x)\,\mathrm{d}x$\qquad
            $\displaystyle \nabla f(x,y)$\par
        }
        \vspace*{1.5cm}
    \end{titlepage}
}

\newcommand{\makebackcover}{%
    \begingroup
    \let\cleardoublepage\clearpage
    \begin{titlepage}
        \thispagestyle{empty}
        \noindent\hspace*{-\parindent}\begin{tikzpicture}[remember picture,overlay]
            \fill[bookblue] (current page.north west) rectangle (current page.south east);
            \fill[booklight,rounded corners=7pt]
                ([xshift=2.2cm,yshift=-3.0cm]current page.north west)
                rectangle
                ([xshift=-2.2cm,yshift=3.0cm]current page.south east);
            \fill[bookteal!80]
                (current page.north west) --
                ([xshift=5.3cm]current page.north west) --
                ([xshift=2.1cm,yshift=-4.4cm]current page.north west) --
                ([yshift=-6.2cm]current page.north west) -- cycle;
            \fill[bookteal!55]
                (current page.south east) --
                ([xshift=-5.8cm]current page.south east) --
                ([xshift=-2.4cm,yshift=4.2cm]current page.south east) --
                ([yshift=6.0cm]current page.south east) -- cycle;
            \foreach \radius in {0.75,1.15,1.55,1.95} {
                \draw[bookteal!55,line width=0.55pt]
                    ([xshift=-2.9cm,yshift=-4.4cm]current page.north east)
                    circle[radius=\radius cm];
            }
            \foreach \radius in {0.65,1.05,1.45} {
                \draw[booklight!65,line width=0.55pt]
                    ([xshift=2.4cm,yshift=3.6cm]current page.south west)
                    circle[radius=\radius cm];
            }
            \fill[bookteal!16,rounded corners=3pt]
                ([xshift=3.5cm,yshift=-7.1cm]current page.north west)
                rectangle
                ([xshift=-5.0cm,yshift=7.2cm]current page.south east);
            \fill[bookteal!30,rounded corners=3pt]
                ([xshift=4.4cm,yshift=-8.0cm]current page.north west)
                rectangle
                ([xshift=-4.1cm,yshift=8.1cm]current page.south east);
            \begin{scope}[shift={([xshift=-0.75cm]current page.center)}]
                \clip[rounded corners=3pt] (-5.25,-4.75) rectangle (5.25,4.75);
                \foreach \v in {-3.6,-3.0,...,3.6} {
                    \draw[bookteal!48,line width=0.45pt,opacity=0.72]
                        plot[domain=-3.6:3.6,samples=45,smooth,variable=\u]
                        ({\u+0.55*\v},{0.28*\u-0.35*\v+0.11*(\u*\u-\v*\v)});
                }
                \foreach \u in {-3.6,-3.0,...,3.6} {
                    \draw[bookblue!42,line width=0.42pt,opacity=0.62]
                        plot[domain=-3.6:3.6,samples=45,smooth,variable=\v]
                        ({\u+0.55*\v},{0.28*\u-0.35*\v+0.11*(\u*\u-\v*\v)});
                }
                \draw[bookteal!70,line width=0.75pt,opacity=0.82]
                    plot[domain=-3.6:3.6,samples=55,smooth,variable=\u]
                    ({\u},{0.28*\u+0.11*\u*\u});
                \draw[bookblue!60,line width=0.7pt,opacity=0.75]
                    plot[domain=-3.6:3.6,samples=55,smooth,variable=\v]
                    ({0.55*\v},{-0.35*\v-0.11*\v*\v});
            \end{scope}
            \draw[bookteal!45,line width=0.65pt]
                ([xshift=3.0cm,yshift=-8.2cm]current page.north west)
                .. controls +(2.2cm,2.0cm) and +(-2.0cm,-2.1cm) ..
                ([xshift=-3.0cm,yshift=9.0cm]current page.south east);
            \draw[bookteal!25,line width=0.55pt]
                ([xshift=3.1cm,yshift=8.0cm]current page.south west)
                .. controls +(2.7cm,-1.7cm) and +(-2.6cm,1.8cm) ..
                ([xshift=-3.1cm,yshift=-8.7cm]current page.north east);
            \node[anchor=west,text=bookteal!72,scale=0.78] at
                ([xshift=2.8cm,yshift=-4.0cm]current page.north west) {%
                $\displaystyle e^{\mathrm{i}\pi}+1=0$};
            \node[anchor=east,text=bookblue!62,scale=0.67] at
                ([xshift=-2.7cm,yshift=4.0cm]current page.south east) {%
                $\displaystyle
                f(x)=\sum_{n=0}^{\infty}\frac{f^{(n)}(a)}{n!}(x-a)^n$};
        \end{tikzpicture}
        \mbox{}
    \end{titlepage}
    \endgroup
}
\fi

% ==================== 新版封面(2026-08 重设计) ====================
% 设计要点:顶部深绿纵向渐变+波浪下缘+两层半透明过渡波(解决深浅绿生硬过渡);
% 深色区螺旋线/同心圆母题;标题+英文副题+作者居中;
% 中部 Fourier 部分和逼近方波曲线;底部浅色波纹(无文字)。
\renewcommand{\maketitle}{%
    \begin{titlepage}
        \thispagestyle{empty}
        \setcounter{page}{0}
        \noindent\hspace*{-\parindent}\begin{tikzpicture}[remember picture,overlay]
            % ---------- 底色 ----------
            \fill[booklight] (current page.north west) rectangle (current page.south east);

            % ---------- 顶部深色区:纵向渐变 + 波浪下缘 ----------
            \shade[top color=bookblue, bottom color=bookteal!55!bookblue]
                (current page.north west) -- (current page.north east) --
                ([yshift=-8.6cm]current page.north east)
                .. controls ([xshift=-5.2cm,yshift=-10.6cm]current page.north east)
                         and ([xshift=6.0cm,yshift=-8.0cm]current page.north west) ..
                ([yshift=-10.0cm]current page.north west) -- cycle;

            % ---------- 过渡层 1:半透明波浪 ----------
            \fill[bookteal!60!bookblue,opacity=0.42]
                ([yshift=-8.6cm]current page.north east)
                .. controls ([xshift=-5.2cm,yshift=-10.6cm]current page.north east)
                         and ([xshift=6.0cm,yshift=-8.0cm]current page.north west) ..
                ([yshift=-10.0cm]current page.north west) --
                ([yshift=-12.6cm]current page.north west)
                .. controls ([xshift=6.5cm,yshift=-11.0cm]current page.north west)
                         and ([xshift=-5.5cm,yshift=-13.2cm]current page.north east) ..
                ([yshift=-11.4cm]current page.north east) -- cycle;

            % ---------- 过渡层 2:更浅的波浪 ----------
            \fill[bookteal!35,opacity=0.35]
                ([yshift=-10.6cm]current page.north east)
                .. controls ([xshift=-6.0cm,yshift=-12.4cm]current page.north east)
                         and ([xshift=5.0cm,yshift=-9.8cm]current page.north west) ..
                ([yshift=-11.8cm]current page.north west) --
                ([yshift=-13.8cm]current page.north west)
                .. controls ([xshift=5.5cm,yshift=-12.6cm]current page.north west)
                         and ([xshift=-6.5cm,yshift=-14.6cm]current page.north east) ..
                ([yshift=-12.8cm]current page.north east) -- cycle;

            % ---------- 深色区装饰:右上螺旋线 ----------
            \begin{scope}[shift={([xshift=-3.4cm,yshift=-4.3cm]current page.north east)}]
                \draw[booklight!45!bookteal,opacity=0.55,line width=0.6pt]
                    plot[domain=0:1440,samples=220,smooth,variable=\t]
                    ({0.0028*\t*cos(\t)},{0.0028*\t*sin(\t)});
            \end{scope}

            % ---------- 深色区装饰:左下同心圆 ----------
            \foreach \r in {0.7,1.25,1.8,2.35,2.9}{
                \draw[booklight!35!bookteal,opacity=0.5,line width=0.5pt]
                    ([xshift=0.4cm,yshift=-7.9cm]current page.north west) circle[radius=\r cm];
            }

            % ---------- 主标题 ----------
            \node[anchor=center,text=booklight] at
                ([yshift=-4.5cm]current page.north)
                {\sffamily\bfseries\fontsize{40}{46}\selectfont \@title};

            % ---------- 分隔短线 ----------
            \draw[booklight!65,line width=1.2pt]
                ([xshift=-2.6cm,yshift=-5.75cm]current page.north) --
                ([xshift=2.6cm,yshift=-5.75cm]current page.north);

            % ---------- 英文副题 ----------
            \node[anchor=center,text=booklight!70!bookteal] at
                ([yshift=-6.55cm]current page.north)
                {\itshape\large Mathematical Analysis \textperiodcentered\
                 A Collection of Exercises};

            % ---------- 作者(紧随标题下方) ----------
            \node[anchor=center,text=booklight!88] at
                ([yshift=-8.2cm]current page.north)
                {\sffamily\large \@author\quad 整理};

            % ---------- 中部:Fourier 部分和逼近方波 ----------
            \begin{scope}[shift={([xshift=-5.2cm,yshift=-18.1cm]current page.north)}]
                % 坐标轴
                \draw[bookgray!55,line width=0.5pt,->] (-0.5,0) -- (10.9,0);
                \draw[bookgray!55,line width=0.5pt,->] (0,-1.75) -- (0,1.95);
                % 目标方波(虚线)
                \draw[bookgray!75,line width=0.6pt,dashed]
                    (0,0.92) -- (5,0.92) (5,-0.92) -- (10,-0.92);
                % N=1
                \draw[bookteal!45,line width=0.7pt]
                    plot[domain=0:360,samples=90,smooth,variable=\x]
                    ({\x/36},{1.18*sin(\x)});
                % N=3
                \draw[bookteal!80,line width=0.8pt]
                    plot[domain=0:360,samples=120,smooth,variable=\x]
                    ({\x/36},{1.18*(sin(\x)+sin(3*\x)/3)});
                % N=7
                \draw[bookblue,line width=0.95pt]
                    plot[domain=0:360,samples=150,smooth,variable=\x]
                    ({\x/36},{1.18*(sin(\x)+sin(3*\x)/3+sin(5*\x)/5+sin(7*\x)/7)});
                % 刻度标注
                \node[anchor=north east,text=bookgray!80,scale=0.62] at (-0.12,0) {$0$};
                \node[anchor=north,text=bookgray!80,scale=0.62] at (5,0) {$\pi$};
                \node[anchor=north,text=bookgray!80,scale=0.62] at (10,0) {$2\pi$};
            \end{scope}
            % 图注
            \node[anchor=center,text=bookgray,scale=0.8] at
                ([yshift=-20.5cm]current page.north)
                {$S_N(x)=\displaystyle\sum_{k=1}^{N}\frac{\sin\bigl((2k-1)x\bigr)}{2k-1}$
                 \ 逐步逼近方波};

            % ---------- 底部浅色波纹(无文字,与顶部波浪呼应) ----------
            \fill[bookteal!22,opacity=0.5]
                ([yshift=2.4cm]current page.south west)
                .. controls ([xshift=6.0cm,yshift=3.9cm]current page.south west)
                         and ([xshift=-6.5cm,yshift=1.0cm]current page.south east) ..
                ([yshift=2.9cm]current page.south east) --
                (current page.south east) -- (current page.south west) -- cycle;
            \fill[bookteal!40!bookblue,opacity=0.28]
                ([yshift=1.3cm]current page.south west)
                .. controls ([xshift=5.5cm,yshift=2.7cm]current page.south west)
                         and ([xshift=-5.5cm,yshift=0.3cm]current page.south east) ..
                ([yshift=1.9cm]current page.south east) --
                (current page.south east) -- (current page.south west) -- cycle;
        \end{tikzpicture}
        \mbox{}
    \end{titlepage}
}

% ==================== 新版封底(纯背景,2026-08 重设计) ====================
% 设计要点:对角渐变底+角部半透明叠层;中下部外摆线(玫瑰形)为视觉主体;
% 底部两层半透明白色波浪;无任何文字。
\newcommand{\makebackcover}{%
    \begingroup
    \let\cleardoublepage\clearpage
    \begin{titlepage}
        \thispagestyle{empty}
        \noindent\hspace*{-\parindent}\begin{tikzpicture}[remember picture,overlay]
            % ---------- 底色:对角渐变 ----------
            \shade[top color=bookblue, bottom color=bookteal!48!bookblue,
                   shading angle=-30]
                (current page.north west) rectangle (current page.south east);

            % ---------- 角部半透明叠层 ----------
            \fill[white,opacity=0.05]
                ([xshift=2.0cm,yshift=1.0cm]current page.north east) circle[radius=7.2cm];
            \fill[white,opacity=0.06]
                ([xshift=-0.5cm,yshift=-1.5cm]current page.north east) circle[radius=3.6cm];

            % ---------- 左下同心圆 ----------
            \foreach \r in {1.0,1.75,2.5,3.25,4.0,4.75}{
                \draw[booklight!40,opacity=0.55,line width=0.55pt]
                    ([xshift=-0.6cm,yshift=-0.4cm]current page.south west) circle[radius=\r cm];
            }
            % ---------- 右下螺旋线 ----------
            \begin{scope}[shift={([xshift=-2.9cm,yshift=3.2cm]current page.south east)}]
                \draw[booklight!45!bookteal,opacity=0.5,line width=0.55pt]
                    plot[domain=0:1260,samples=200,smooth,variable=\t]
                    ({0.0036*\t*cos(\t)},{0.0036*\t*sin(\t)});
            \end{scope}
            % ---------- 左上斜线组 ----------
            \foreach \i in {0,1,2}{
                \draw[bookteal!60!booklight,opacity=0.5,line width=0.6pt]
                    ([xshift={1.6cm+\i*0.45cm}]current page.north west) --
                    ([yshift={-1.6cm-\i*0.45cm}]current page.north west);
            }

            % ---------- 中下部视觉主体:外摆线(玫瑰形) ----------
            \begin{scope}[shift={([yshift=-16.6cm]current page.north)}]
                % 外圈
                \draw[booklight!35,opacity=0.6,line width=0.6pt]
                    (0,0) circle[radius=4.9cm];
                \draw[booklight!22,opacity=0.5,line width=0.5pt]
                    (0,0) circle[radius=5.25cm];
                % 主玫瑰线
                \draw[booklight!55!bookteal,opacity=0.8,line width=0.85pt]
                    plot[domain=0:1080,samples=400,smooth,variable=\t]
                    ({0.335*(8*cos(\t)-5*cos(2.6667*\t))},
                     {0.335*(8*sin(\t)-5*sin(2.6667*\t))});
                % 内层玫瑰线
                \draw[booklight!38!bookteal,opacity=0.6,line width=0.6pt]
                    plot[domain=0:1080,samples=400,smooth,variable=\t]
                    ({0.21*(8*cos(\t)-5*cos(2.6667*\t))},
                     {0.21*(8*sin(\t)-5*sin(2.6667*\t))});
            \end{scope}

            % ---------- 底部半透明波浪 ----------
            \fill[white,opacity=0.06]
                ([yshift=4.6cm]current page.south west)
                .. controls ([xshift=6.5cm,yshift=6.4cm]current page.south west)
                         and ([xshift=-6.0cm,yshift=2.6cm]current page.south east) ..
                ([yshift=5.2cm]current page.south east) --
                (current page.south east) -- (current page.south west) -- cycle;
            \fill[booklight,opacity=0.08]
                ([yshift=2.6cm]current page.south west)
                .. controls ([xshift=5.5cm,yshift=4.3cm]current page.south west)
                         and ([xshift=-6.5cm,yshift=1.0cm]current page.south east) ..
                ([yshift=3.3cm]current page.south east) --
                (current page.south east) -- (current page.south west) -- cycle;
        \end{tikzpicture}
        \mbox{}
    \end{titlepage}
    \endgroup
}
\makeatother

% 页脚右下方常态化水印(署名 + 仓库地址,加粗黑色),已在 main/plain 分页样式中定义,所有页面均生效.

\begin{document}
\mainmatter
\pagestyle{main}
\setcounter{chapter}{2}
\def\BookEnd{\end{document}}
\fi

\chapter{函数}

\begin{ti}
【题1】用覆盖定理证明零点存在性定理
\end{ti}

\textbf{证明}:给定连续函数$f,f(a)<0,f(b)>0,$ \\

假设函数$f\text{在}[a,b]$上无零点,则对$\forall\,x_0\in[a,b],$由连续函数保号性, \\

$\exists\,\delta_x>0,\;\text{s.t.}\,\forall\,x\in O_{\delta_x}(x_0)\cap[a,b],\text{都有}f(x)f(x_0)>0,$ \\

又$\displaystyle\bigcup_{x\in[a,b]}O_{\delta_x}(x)\text{构成}[a,b]$的一个开覆盖,那么,由$Lebesgue$定理, \\

$\exists\,\zeta>0,\;\text{s.t.}\,x',x''\in[a,b]\text{且}\left\lvert x'-x''\right\rvert<\zeta\text{时},x',x''\text{可被包在同一个}O_{\delta}(x)$里, \\

在$[a,b]\text{中插入一系列点,连同端点一起,记为:}a=x_0<x_1<x_2<...<x_n=b,$ \\

$\text{且满足}\left\lvert x_i-x_{i-1}\right\rvert<\zeta,\,i=1,2,...,n,$ \\

由插入点的特点可知,$f(x_i)f(x_{i-1})>0\Longrightarrow f(a)f(b)>0,$矛盾.

\begin{ti}
【题2】设$f\in C[a,b),\{x_n\},\{y_n\}\subset[a,b),\,\,\text{满足}:$\\

$\displaystyle\lim_{n\to\infty}x_n=\lim_{n\to\infty}y_n=b,\displaystyle\lim_{n\to\infty}f(x_n)=A<B=\lim_{n\to\infty}f(y_n)$ \\

证明:$\forall\,\eta\in(A,B),\;\exists\,\{z_n\}\subset[a,b),\;\text{s.t.}\,\displaystyle\lim_{n\to\infty}z_n=b,\lim_{n\to\infty}f(z_n)=\eta$
\end{ti}

\textbf{证明}:取$\varepsilon=\min\{\eta-A,B-\eta\},\text{则}\exists\,N,$ \\

$n>N\text{时},\,f(x_n)<A+\varepsilon=\eta,f(y_n)>B-\varepsilon=\eta,$ \\

又$x_n,y_n\in[a,b),f\in C[a,b)$,由零点存在性定理, \\

$\exists\,z_n\text{介于}x_n,y_n\text{之间},\text{且}f(z_n)=\eta,$ \\

又由夹逼定理,可得$\displaystyle\lim_{n\to\infty}z_n=b.$

\begin{ti}
【题3】设$f$在$(a,b)$上一致连续,证明:存在两个有限的单侧极限$f(a^+),f(b^-)$
\end{ti}

\textbf{证明}:$\forall\,\varepsilon>0,\;\exists\,\delta>0,\;\text{s.t.}\,\left\lvert x'-x''\right\rvert<\delta\text{且}x',x''\in(a,b)\text{时},\left\lvert f(x')-f(x'')\right\rvert<\varepsilon,$ \\

$\therefore\,x_1,x_2\in(a,a+\delta)\text{时},\left\lvert f(x_1)-f(x_2)\right\rvert<\varepsilon,$ \\

由$Cauchy$收敛准则,$\displaystyle\lim_{x\to a^+}f(x)$存在且有限.同理$f(b^-)$有限,得证.

\begin{ti}
【题4】证明:$f(x)=x\sin{\dfrac{1}{x}}$在$(0,+\infty)$上一致连续
\end{ti}

\textbf{证明}:$\left\lvert f(x)-f(y)\right\rvert=\left\lvert x\sin{\dfrac{1}{x}}-y\sin{\dfrac{1}{y}}\right\rvert=\left\lvert x\sin{\dfrac{1}{x}}-y\sin{\dfrac{1}{x}}+y\sin{\dfrac{1}{x}}-y\sin{\dfrac{1}{y}}\right\rvert$ \\

$<\left\lvert \sin{\dfrac{1}{x}}\right\rvert\left\lvert x-y\right\rvert+y\left\lvert \sin{\dfrac{1}{x}-\dfrac{1}{y}}\right\rvert<\left\lvert x-y\right\rvert+y\left\lvert \sin{\dfrac{1}{x}-\dfrac{1}{y}}\right\rvert,$ \\

利用$\sin{x}\text{是}Lipschitz\text{函数},\,\left\lvert \sin{\dfrac{1}{x}-\dfrac{1}{y}}\right\rvert<\left\lvert \dfrac{1}{x}-\dfrac{1}{y}\right\rvert,$ \\

$\therefore\left\lvert f(x)-f(y)\right\rvert<\left\lvert x-y\right\rvert+y\left\lvert \dfrac{1}{x}-\dfrac{1}{y}\right\rvert=\left\lvert x-y\right\rvert+\dfrac{1}{x}\left\lvert x-y\right\rvert,$ \\

同理$\left\lvert f(x)-f(y)\right\rvert<\left\lvert x-y\right\rvert+\dfrac{1}{y}\left\lvert x-y\right\rvert,$ \\

又$\left\lvert f(x)-f(y)\right\rvert<\left\lvert x\sin{\dfrac{1}{x}}\right\rvert+\left\lvert y\sin{\dfrac{1}{y}}\right\rvert<x+y,$ \\

对$\forall\,\varepsilon>0,\;\text{取}\delta=\varepsilon(1+\dfrac{2}{\varepsilon})^{-1},$我们断言:$\left\lvert x-y\right\rvert<\delta\text{时},\left\lvert f(x)-f(y)\right\rvert<\varepsilon,$ \\

分情况讨论, \\

\romannum{1} 若$x<\dfrac{\varepsilon}{2},\,y<\dfrac{\varepsilon}{2},\text{则}\left\lvert f(x)-f(y)\right\rvert<x+y<\varepsilon,$ \\

\romannum{2} 若$x\geq\dfrac{\varepsilon}{2}(y\geq\dfrac{\varepsilon}{2}\text{时类似可这样处理}),\text{则}\left\lvert f(x)-f(y)\right\rvert<\left\lvert x-y\right\rvert+\dfrac{1}{x}\left\lvert x-y\right\rvert$ \\

$\leq\left\lvert x-y\right\rvert+\dfrac{2}{\varepsilon}\left\lvert x-y\right\rvert<(1+\dfrac{2}{\varepsilon})\delta=\varepsilon,$ \\

即$\forall\,\varepsilon>0,\;\exists\,\delta=\varepsilon(1+\dfrac{2}{\varepsilon})^{-1},\left\lvert x-y\right\rvert<\delta\text{时},\left\lvert f(x)-f(y)\right\rvert<\varepsilon,$ \\

即$f(x)=x\sin{\dfrac{1}{x}}\text{在}(0,+\infty)$上一致连续. \\

\textbf{注}:本题中$\left\lvert f(x)-f(y)\right\rvert=\left\lvert x\sin{\dfrac{1}{x}}-y\sin{\dfrac{1}{x}}+y\sin{\dfrac{1}{x}}-y\sin{\dfrac{1}{y}}\right\rvert$的处理方法在解决这种差值估计当中十分常见,
但是本题中这种方法在0附近失效了,不能直接解决问题,所以要分段去用不同的放缩,在0附近利用$f$趋近于0来放缩.

\begin{ti}
【题5】设$g\text{是区间}I\text{上的单调函数}$,证明:值域$g(I)$为区间的充要条件为$g\in C^0(I)$
\end{ti}

\textbf{证明}:充分性易得,下证必要性,先设$\displaystyle g$是单调递增的, \\

若$\displaystyle g$在内点$\displaystyle x_0\in I$处不连续,则$\displaystyle g(x_0^-)<g(x_0^+)$,取$\displaystyle y\in(g(x_0^-),g(x_0^+))$且$\displaystyle y\neq g(x_0)$, \\

由单调性可知$\displaystyle y\notin g(I)$,与$\displaystyle g(I)$是区间矛盾, \\

若$\displaystyle I$含左端点$\displaystyle a$,且$\displaystyle g$在$\displaystyle a$处相对右不连续,则$\displaystyle g(a)<g(a^+)$, \\

任取$\displaystyle y\in(g(a),g(a^+))$,对$\displaystyle x=a$有$\displaystyle g(x)<y$,对$\displaystyle x>a$有$\displaystyle g(x)\geq g(a^+)>y$, \\

故$\displaystyle y\notin g(I)$,矛盾, \\

若$\displaystyle I$含右端点$\displaystyle b$,且$\displaystyle g$在$\displaystyle b$处相对左不连续,则$\displaystyle g(b^-)<g(b)$,\\

任取$\displaystyle y\in(g(b^-),g(b))$,对$\displaystyle x<b$有$\displaystyle g(x)\leq g(b^-)<y$,而$\displaystyle g(b)>y$,故$\displaystyle y\notin g(I)$,矛盾, \\

所以单调递增时$\displaystyle g\in C^0(I)$.若$\displaystyle g$单调递减,对$\displaystyle-g$应用上述结论即可,故$\displaystyle g\in C^0(I).$

\begin{ti}
【题6】设$f$是$(-\infty,+\infty)$上的单调函数,在每一点定义$g(x)=f(x^+)$,证明:$g\text{是}(-\infty,+\infty)$上处处右连续的函数
\end{ti}

\textbf{证明}:即证$\forall\,x_0\in\mathbb{R},\;\displaystyle\lim_{x\to x_0^+}f(x^+)=f(x_0^+),$ \\

不妨设$f$是单调递增函数,则$x>x_0\text{时},f(x_0^+)\leq f(x)\leq f(x^+),$ \\

又$f(x_0^+)=\displaystyle\inf_{x>x_0}f(x),$由确界的刻画可知, \\

$\forall\,\varepsilon>0,\;\exists\,u_0>x_0,\;\text{s.t.}\,f(x_0^+)+\varepsilon>f(u_0),$ \\

令$\delta=u_0-x_0>0,\,x\in(x_0,x_0+\delta)\text{时},f(x^+)\leq f(u_0)<f(x_0^+)+\varepsilon,$ \\

$\therefore\,\forall\varepsilon>0,\;\exists\,\delta>0,x\in(x_0,x_0+\delta)\text{时},\left\lvert f(x^+)-f(x_0^+)\right\rvert<\varepsilon,$ \\

$\therefore\displaystyle\lim_{x\to x_0^+}f(x^+)=f(x_0^+)$对$\forall\,x_0\in\mathbb{R}$成立,即$g(x)$在$(-\infty,+\infty)$上处处右连续.

\begin{ti}
【题7】设函数$f$在区间$I$内只有可去间断点,定义$g(x)=\displaystyle\lim_{t\to x}f(t),x\in I,$证明:$g\in C(I)$
\end{ti}

\textbf{证明}:因为$\displaystyle f$只有可去间断点,所以$\displaystyle g$在$\displaystyle I$中每一点都有定义, \\

给定$\displaystyle\varepsilon>0$,由极限定义,存在$\displaystyle\delta>0$,当$\displaystyle0<|t-x_0|<\delta$时,$\displaystyle\left\lvert f(t)-g(x_0)\right\rvert<\dfrac{\varepsilon}{2},$ \\

若$\displaystyle|x-x_0|<\dfrac{\delta}{2}$,令$\displaystyle y\to x$且$\displaystyle0<|y-x|<\dfrac{\delta}{2}-|x-x_0|$,则$\displaystyle|y-x_0|<\delta$,从而 \\

$\displaystyle g(x_0)-\dfrac{\varepsilon}{2}\leq\lim_{y\to x}f(y)=g(x)\leq g(x_0)+\dfrac{\varepsilon}{2},$ 故$\displaystyle\left\lvert g(x)-g(x_0)\right\rvert\leq\dfrac{\varepsilon}{2}<\varepsilon$,\\

因此$\displaystyle\lim_{x\to x_0}g(x)=g(x_0)$对任意$\displaystyle x_0\in I$成立,即$\displaystyle g\in C(I).$

\begin{ti}
【题8】证明:若$f$在$[0,+\infty)$上连续且有界,则对每个给定的$\lambda$,存在一个数列$\{x_n\}$满足要求: \\
$\displaystyle(1)\lim_{n\to\infty}x_n=+\infty;\quad(2)\lim_{n\to\infty}[f(x_n+\lambda)-f(x_n)]=0$
\end{ti}

\textbf{证明}:若$\lambda=0,\text{取}x_n=n$即可,下设$\lambda\neq 0$, \\

考虑函数$\displaystyle g(x)=f(x+\lambda)-f(x)$,则$\displaystyle g$在$\displaystyle[\max\{0,-\lambda\},+\infty)$上连续且有界, \\

下面证明:$\displaystyle\forall\,\varepsilon>0,\forall\,M>0,\;\exists\,x\geq M,\;\text{s.t.}\,\left\lvert g(x)\right\rvert<\varepsilon,$ \\

用反证法,假设$\displaystyle\exists\,\varepsilon_0>0,M_0>0,\;\text{s.t.}\,\forall\,x\geq M_0,\left\lvert g(x)\right\rvert\geq\varepsilon_0$, \\

令$\displaystyle L=\max\{0,-\lambda,M_0\}$,因为$\displaystyle g$在连通区间$\displaystyle[L,+\infty)$上连续且不取零值, \\

所以由介值定理,$\displaystyle g$在该区间上恒取同一符号, \\

即$\displaystyle g(x)\geq\varepsilon_0$恒成立或$\displaystyle g(x)\leq-\varepsilon_0$恒成立, \\

下设前一种情形,后一种情形同理, \\

当$\displaystyle\lambda>0$时,$\displaystyle f(L+\lambda)-f(L)\geq\varepsilon_0,$ \\

\phantom{当$\displaystyle\lambda>0$时,$\displaystyle\;$}$\displaystyle f(L+2\lambda)-f(L+\lambda)\geq\varepsilon_0,$ \\
\phantom{当$\displaystyle\lambda>0$时,$\displaystyle\;f(L+2\lambda)$}\vdots \\
\phantom{当$\displaystyle\lambda>0$时,$\displaystyle\;$}$\displaystyle f(L+n\lambda)-f(L+n\lambda-\lambda)\geq\varepsilon_0,$ \\

\phantom{当$\displaystyle\lambda>0\;$}$\displaystyle\therefore\,f(L+n\lambda)-f(L)\geq n\varepsilon_0,$ \\

\phantom{当$\displaystyle\lambda>0\;$}$\displaystyle\therefore\lim_{n\to\infty}f(L+n\lambda)=+\infty$,与$\displaystyle f(x)$在$\displaystyle[0,+\infty)$上有界矛盾, \\

当$\displaystyle\lambda<0$时,$\displaystyle f(L)-f(L-\lambda)\geq\varepsilon_0,$ \\

\phantom{当$\displaystyle\lambda<0$时,$\displaystyle\;$}$\displaystyle f(L-\lambda)-f(L-2\lambda)\geq\varepsilon_0,$ \\
\phantom{当$\displaystyle\lambda<0$时,$\displaystyle\;f(L-\lambda)$}\vdots \\
\phantom{当$\displaystyle\lambda<0$时,$\displaystyle\;$}$\displaystyle f(L-n\lambda+\lambda)-f(L-n\lambda)\geq\varepsilon_0,$ \\

\phantom{当$\displaystyle\lambda<0\;$}$\displaystyle\therefore\,f(L)-f(L-n\lambda)\geq n\varepsilon_0,$ \\

\phantom{当$\displaystyle\lambda<0\;$}$\displaystyle\therefore\lim_{n\to\infty}f(L-n\lambda)=-\infty$,与$\displaystyle f(x)$在$\displaystyle[0,+\infty)$上有界矛盾, \\

综上,$\forall\,\varepsilon>0,\forall\,M>0,\;\exists\,x\geq M,\;\text{s.t.}\,\left\lvert g(x)\right\rvert<\varepsilon,$ \\

取$\varepsilon=1,M=1,\;\exists\,x_1\geq 1,\text{且}\left\lvert g(x_1)\right\rvert<1,$ \\

取$\varepsilon=\dfrac{1}{2},M=2,\;\exists\,x_2\geq 2,\text{且}\left\lvert g(x_2)\right\rvert<\dfrac{1}{2},$ \\
\phantom{取$\varepsilon=1,M=1,$}\vdots \\
取$\varepsilon=\dfrac{1}{n},M=n,\;\exists\,x_n\geq n,\text{且}\left\lvert g(x_n)\right\rvert<\dfrac{1}{n},$ \\

$\therefore\displaystyle\lim_{n\to\infty}x_n=+\infty,\lim_{n\to\infty}\left\lvert f(x_n+\lambda)-f(x_n)\right\rvert=0,$ \\

综上,$\forall\,\lambda\in\mathbb{R},\;\exists\{x_n\},\;\text{s.t.}\,\displaystyle\lim_{n\to\infty}x_n=+\infty,\lim_{n\to\infty}[f(x_n+\lambda)-f(x_n)]=0.$

\begin{ti}
【题9】证明:不等于常数的连续周期函数一定有最小正周期
\end{ti}

\textbf{证明}:记集合$T=\left\{t>0\,|\;f(x+t)=f(x)\right\}$,函数$f(x)$为不等于常数的连续周期函数, \\

$\because\,f(x)$是周期函数,$\quad\therefore\,T$非空,又$t>0,\quad\therefore\,T$有下界, \\

由确界原理,$\exists\,T_0=\inf{T},T_0\geq 0$,下证:$T_0>0\text{且}T_0\in T,$ \\

\ding{172} 若$\displaystyle T_0=0$,则对每个$\displaystyle n\in\mathbb{N}^+$,可取$\displaystyle t_n\in T$使$\displaystyle0<t_n<\dfrac{1}{n}$, \\

任取$\displaystyle x_1<x_2$,令$\displaystyle k_n$为满足$\displaystyle k_nt_n\geq x_2-x_1$的最小正整数,并令$\displaystyle r_n=x_2-x_1-(k_n-1)t_n$, \\

则$\displaystyle0<r_n\leq t_n<\dfrac{1}{n}$,故$\displaystyle r_n\to0^+$,又$\displaystyle(k_n-1)t_n$为$\displaystyle f$的周期, \\

$\therefore\displaystyle f(x_2)=f(x_1+(k_n-1)t_n+r_n)=f(x_1+r_n)\to f(x_1),$ \\

故$\displaystyle f(x_2)=f(x_1)$.由$\displaystyle x_1,x_2$的任意性,$\displaystyle f$为常值函数,与题目条件矛盾!$\displaystyle\quad\therefore\,T_0>0,$ \\

\ding{173} $\because\,T_0=\inf{T},\;\therefore\exists\,\{t_n\}\subseteq T,\;\text{s.t.}\,\displaystyle\lim_{n\to\infty}t_n=T_0,$ \\

又$f(x)=f(x+t_n)\text{对}\forall\,n\in\mathbb{N}^+$成立,$\quad\therefore\,f(x)=\displaystyle\lim_{n\to\infty}f(x+t_n),$ \\

由$Heine$定理,$\displaystyle\lim_{n\to\infty}f(x+t_n)=f(x+T_0),\quad\therefore\,f(x)=f(x+T_0),\text{即}T_0\in T,$ \\

综上,$T_0$为$f(x)$的最小正周期,即$f(x)$有最小正周期. \\

\textbf{注记}:关键在于,当周期很小的时候,可以用整数倍周期很好的逼近任意实数.

\begin{ti}
【题10】设$f$在区间$[0,n]$连续,$f(0)=f(n),n\in\mathbb{N}^+$,证明:至少有$n$对不同的$(x,y)$使得$f(x)=f(y)$,同时$x-y$为非$0$整数
\end{ti}

\textbf{证明}:对$n$进行归纳,$n=1$时显然成立, \\

假设$n=k-1$时结论成立,$\;n=k$时,设$g(x)=f(x+1)-f(x),$ \\

则$g(0)+g(1)+...+g(n-1)=0,\;\therefore\,g(0),g(1),...g(n-1)$不能恒正或恒负, \\

$\therefore\exists\,k_1,k_2\in\{0,1,...,n-1\},\;\text{s.t.}\,g(k_1)\geq 0,g(k_2)\leq 0,k_1\neq k_2,$ \\

$\therefore\exists\,x_0\text{介于}k_1,k_2\text{之间},\;\text{s.t.}\,g(x_0)=0,$ \\

记$\displaystyle\phi(x)=\begin{cases}
    \raisebox{0.45ex}[0pt][0pt]{$\displaystyle x$}&\raisebox{0.45ex}[0pt][0pt]{$\displaystyle\quad x\in[0,x_0]$} \\
    \displaystyle x+1&\displaystyle\quad x\in(x_0,n-1]
\end{cases}$,并令$\displaystyle h(x)=f(\phi(x))$,
\\

因为$\displaystyle f(x_0)=f(x_0+1)$,所以$\displaystyle h\in C[0,n-1]$, \\

由归纳假设,$\displaystyle h$至少有$\displaystyle n-1$对不同的$\displaystyle(x_i,y_i)$,满足$\displaystyle h(x_i)=h(y_i)$且$\displaystyle x_i-y_i\in\mathbb{Z}\setminus\{0\}$, \\

易得函数$\displaystyle\phi$在$\displaystyle[0,n-1]$上严格递增,故为单射, \\

对每一对$\displaystyle(x_i,y_i)$,有$\displaystyle f(\phi(x_i))=f(\phi(y_i))$, \\

且$\displaystyle\phi(x_i)-\phi(y_i)$等于$\displaystyle x_i-y_i$加上$\displaystyle-1,0,1$中的一个,因而仍为整数; \\

又由单射性可知该整数非零, 单射性还保证不同点对的像仍不相同, \\

又$\displaystyle\phi([0,n-1])=[0,x_0]\cup(x_0+1,n]$,所以$\displaystyle x_0+1$不在像集中, \\

从而新增点对$\displaystyle(x_0,x_0+1)$不与上述$\displaystyle n-1$对重复, \\

因此$\displaystyle f$至少有$\displaystyle n$对不同的$\displaystyle(x,y)$满足$\displaystyle f(x)=f(y)$且$\displaystyle x-y$为非零整数,\\

由数学归纳法可知,至少有$\displaystyle n$对不同的$\displaystyle(x,y)$使得$\displaystyle f(x)=f(y)$,同时$\displaystyle x-y$为非$\displaystyle0$整数.



\BookEnd
